\documentclass{article}
\usepackage[utf8]{inputenc}
\usepackage{amsmath}
\usepackage{geometry}
\geometry{margin=2cm}
\begin{document}

\section*{Questão 97 — ENEM 2024 — Ciências da Natureza}

\subsection*{Interpretação e Estratégia}
A questão apresenta três tubos de ensaio com culturas bacterianas cultivadas no mesmo meio nutritivo e em condições físico-químicas ideais. A distribuição das bactérias, representadas por pontos vermelhos, varia em cada tubo, refletindo o tipo de metabolismo energético dos microrganismos em relação à disponibilidade de oxigênio.

O enunciado solicita a classificação dos microrganismos dos tubos 1, 2 e 3 como aeróbio estrito, anaeróbio estrito ou anaeróbio facultativo.

Para isso, a estratégia é analisar para cada tubo onde há maior concentração bacteriana:
\begin{itemize}
  \item Topo do tubo – maior concentração de oxigênio dissolvido.
  \item Fundo do tubo – ambiente anóxico, ausência de oxigênio.
  \item Distribuição uniforme – capacidade de crescer com ou sem oxigênio.
\end{itemize}

\subsection*{Conteúdo e Revisão}
\begin{table}[h!]
  \centering
  \begin{tabular}{|p{5cm}|p{5cm}|p{5cm}|}
    \hline
    Conceito & Ideia-chave & Aplicação \\
    \hline
    Metabolismo bacteriano & Bactérias podem ser aeróbias (necessitam oxigênio), anaeróbias estritas (não toleram oxigênio) e anaeróbias facultativas (crescem com/sem oxigênio) & Posição das bactérias no tubo indica metabolismo \\
    \hline
    Oxigênio dissolvido em meio líquido & Oxigênio é mais concentrado na superfície do meio e praticamente ausente no fundo & Explica os locais de crescimento das bactérias em tubos \\
    \hline
  \end{tabular}
  \caption{Revisão dos conceitos-chave para a questão}
\end{table}

A teoria baseia-se no fato de que bactérias aeróbias estritas respiram com oxigênio e crescem no topo do meio líquido, anaeróbias estritas evitam oxigênio e crescem no fundo do tubo, e facultativas adaptam-se aos dois ambientes, colonizando todo o meio.

\subsection*{Resolução e "Pulo do Gato"}
Observa-se:
\begin{itemize}
  \item No tubo 1, bactérias estão no topo, indicando metabolismo aeróbio estrito.
  \item No tubo 2, bactérias no fundo, caracterizando anaeróbio estrito.
  \item No tubo 3, bactérias distribuídas uniformemente, indicando anaeróbio facultativo.
\end{itemize}

\textbf{Conclusão:} a sequência correta é \textbf{aeróbio estrito, anaeróbio estrito, anaeróbio facultativo} (alternativa C).

\textbf{Pulo do gato:} "Onde há oxigênio, só aeróbios crescem; onde não há, anaeróbios estritos; quem cresce em tudo é o facultativo." 

\subsection*{Armadilhas e Distratores}
\begin{itemize}
  \item Confundir os anaeróbios facultativos com estritos ao interpretar a distribuição das bactérias.
  \item Trocar o significado do crescimento na superfície do tubo, atribuindo erradamente a anaeróbios facultativos.
  \item Inverter a relação entre oxigênio disponível e local de crescimento bacteriano.
  \item Não perceber que anaeróbios estritos não crescem na presença de oxigênio.
\end{itemize}

\subsection*{Código da Questão}
\texttt{CN_MICROBIOLOGIA_CIENCIAS_001}

\vspace{1cm}
\begin{center}
\textbf{\large Alternativa Correta: C}
\end{center}

\vspace{0.5cm}
\noindent
\textit{Resolução oficial:} 

A distribuição das bactérias nos tubos indica suas preferências metabólicas quanto ao oxigênio: no tubo 1, concentração na superfície revela aeróbio estrito; no tubo 2, concentração no fundo indica anaeróbio estrito; no tubo 3, distribuição generalizada mostra anaeróbio facultativo. Portanto, a alternativa que classifica corretamente é \textbf{C}.

\end{document}